\scanpage{267}
\mspage{781}

\unsure{mant|me} by the commutativity of the diagrams
\[
(107)\quad
\diagram{\begin{tikzcd}[ampersand replacement=\&, column sep=large,row sep=small]
H(A) \arrow[dd,"(100)"',"\sim"] \arrow[rd,"\substack{i_A\\\sim}"] \& \\
\& H(\unsure{\mathbb{S}}_x^{\sim},\unsure{\mathbb{D}}_x^{*\sim}) \\
H(A') \arrow[ru,"\sim","i_{A'}"'] \&
\end{tikzcd}}
\;.
\]

To verify the fact that the functor $\underline{\mathrm{Gr}}(H)$ of (102)
is \unsure{essentially} \unsure{constant},
I consider the two functors (independent of the choice of
$\unsure{\mathbb{S}}_x^{\sim},\unsure{\mathbb{D}}_x^{*\sim}$)
\[
(108)\quad
\left\{
\diagram{\begin{tikzcd}[ampersand replacement=\&, column sep=small,row sep=large]
A \arrow[r,mapsto] \& \unsure{\mathbb{S}}_{xA}^{\sim},
\& A \arrow[r,mapsto] \& \unsure{\mathbb{D}}_{x,A}^{*\sim} \\
\underline{\mathcal{A}}_x \arrow[r] \& \underline{\unsure{\mathrm{Revuni}}}(\unsure{\mathbb{S}}_x)
\& \underline{\mathcal{A}}_x \arrow[r] \& \underline{\unsure{\mathrm{Revuni}}}(\unsure{\mathbb{D}}_x^*)
\end{tikzcd}}
\right.
\;,
\]
which factor into \uline{equivalences}
\[
(109)\quad
\diagram{\begin{tikzcd}[ampersand replacement=\&, column sep=small]
\underline{\mathrm{Gr}}(\underline{\mathcal{A}}_x)
\arrow[r,"\substack{F\\\approx}"] \&
\underline{\unsure{\mathrm{Revuni}}}(\unsure{\mathbb{S}}_x),
\& \underline{\mathrm{Gr}}(\underline{\mathcal{A}}_x)
\arrow[r,"\substack{G\\\approx}"] \&
\underline{\unsure{\mathrm{Revuni}}}(\unsure{\mathbb{D}}_x^*)
\end{tikzcd}}
\]
This said, \inserted{using
$\unsure{\mathbb{S}}_x^{\sim},\unsure{\mathbb{D}}_x^{*\sim}$,}
we have for $A\in\mathrm{Ob}\,\underline{\mathcal{A}}_x=\mathcal{A}_x$ also,
isomorphisms functorial in $A$
\[
(110)\quad\left\{
\begin{array}{l}
\unsure{\mathbb{S}^{\sim}|H}(A)
\simeq\mathrm{Isom}_{\underline{\unsure{\mathrm{Revuni}}}(\unsure{\mathbb{S}}_x)}
\deleted{\unsure{\mathbb{S}}}
\left(\underbrace{F(A)}_{\substack{\Vert\\\unsure{\mathbb{S}}_{x,A}^{\sim}}},\unsure{\mathbb{S}}^{\sim}\right)
\\[1.8em]
\unsure{\mathbb{D}}_x^{*\sim}(\unsure{\mathrm{Ep}|E_p}A)
\simeq\mathrm{Isom}_{\underline{\unsure{\mathrm{Revuni}}}(\unsure{\mathbb{D}}_x^*)}
\left(\underbrace{G(A)}_{\substack{\Vert\\\unsure{\mathbb{D}}_{x,A}^{*\sim}}},\unsure{\mathbb{D}}_x^{*\sim}\right),
\end{array}\right.
\]
and by the universal property of $\underline{\mathrm{Gr}}(\underline{\mathcal{A}}_x)$,
these functorial isomorphisms extend over the category
$\underline{\mathrm{Gr}}(\underline{\mathcal{A}}_x)$.
Thus the situation is one where we have a commutative group $\unsure{T|\overline T}$ $(=T_x)$,
three connected $T$-groupoids $\mathcal{G},\mathcal{G}',\mathcal{G}''$

\margin{266}

\margin{\begin{tabular}{c}
266\\
\texttt{https://grothendieck.umontpellier.fr}\\
Université de Montpellier, UMR CNRS 5149, Case courrier 051, Place Eugène Bataillon\\
34095 Montpellier cedex 5 - France
\end{tabular}}

\scanpage{268}

(namely $\underline{\mathrm{Gr}}(\underline{\mathcal{A}}_x)$,
$\underline{\unsure{\mathrm{Revuni}}}(\unsure{\mathbb{S}}_x)$,
$\underline{\unsure{\mathrm{Revuni}}}(\unsure{\mathbb{D}}_x^*)$),
equivalences \deleted{\unsure{of}} of $T$-groupoids
\[
(111)\quad
\diagram{\begin{tikzcd}[ampersand replacement=\&, column sep=large]
\mathcal{G} \arrow[r,"\substack{F\\\approx}"] \& \mathcal{G}',
\& \mathcal{G} \arrow[r,"\substack{G\\\approx}"] \& \mathcal{G}''
\end{tikzcd}}
\]
and objects $x'\in\mathrm{Ob}\,\mathcal{G}'$,
$x''\in\mathrm{Ob}\,\mathcal{G}''$ (namely
$\unsure{\mathbb{S}}^{\sim}$ and $\unsure{\mathbb{D}}_x^{*\sim}$),
and if we consider the functors
\[
(112)\quad\left\{
\diagram{\begin{tikzcd}[ampersand replacement=\&, column sep=small,row sep=large]
\mathcal{G} \arrow[r,"f=f_{x'}"] \& \underline{\mathrm{Tors}}(T),
\& \mathcal{G} \arrow[r,"g=g_{x''}"] \& \underline{\mathrm{Tors}}(T) \\
x \arrow[r,mapsto] \& \mathrm{Hom}_{\mathcal{G}'}(F(x),x'),
\& x \arrow[r,mapsto] \& \mathrm{Hom}_{\mathcal{G}''}(G(x),x'')
\end{tikzcd}}
\right.
\]
and the corresponding functor
\[
(113)\quad\left\{
\diagram{\begin{tikzcd}[ampersand replacement=\&, column sep=large,row sep=large]
\mathcal{G} \arrow[r] \& \underline{\mathrm{Tors}}(T) \\
x \arrow[r,mapsto] \&
\mathrm{Hom}_{\underline{\mathrm{Tors}}(\unsure{T|G})}(f(x),g(x))
\end{tikzcd}}
\right.
\;,
\]
to prove that this last functor is essentially constant.
We have then, if \deleted{$F',G'$ denote}
\[
(114)\quad
\diagram{\begin{tikzcd}[ampersand replacement=\&, column sep=large]
F_1\deleted{\unsure{f}}:\mathcal{G}' \arrow[r,"\approx"] \& \mathcal{G},
\& \unsure{{}'G_1|G_1}:\mathcal{G}'' \arrow[r,"\approx"] \& \mathcal{G}
\end{tikzcd}}
\]
denote quasi-inverses of $F,G$, isomorphisms functorial in $x$
\[
(115)\quad\left\{
\begin{array}{l}
f(x)\overset{\mathrm{def}}{=}\mathrm{Hom}_{\mathcal{G}'}(F(x),x')
\simeq\mathrm{Hom}_{\mathcal{G}}(x,F_1(x')) \\
g(x)\overset{\mathrm{def}}{=}\mathrm{Hom}_{\mathcal{G}''}(G(x),x'')
\simeq\mathrm{Hom}_{\mathcal{G}}(x,F_2(x''))
\end{array}\right.
\]
whence \unsure{follows|follow} isomorphisms of $T$-torsors
\[
(116)\quad
\unsure{\mathrm{Isom}}(f(x),g(x))
\simeq\unsure{\mathrm{Isom}}(F_1(x'),F_2(x''))
\]
whence the assertion.

\margin{267}

\margin{\begin{tabular}{c}
267\\
\texttt{https://grothendieck.umontpellier.fr}\\
Université de Montpellier, UMR CNRS 5149, Case courrier 051, Place Eugène Bataillon\\
34095 Montpellier cedex 5 - France
\end{tabular}}

\scanpage{269}
\mspage{782}

\uline{Definition} We say that we have a \uline{correlation}
between a universal covering $\unsure{\mathbb{S}}_x^{\sim}$ of
$\unsure{\mathbb{S}}_x$, and another $\unsure{\mathbb{D}}_x^{*\sim}$
of $\unsure{\mathbb{D}}_x^*$, when we have given ourselves an element of
$H(\unsure{\mathbb{S}}_x^{\sim},\unsure{\mathbb{D}}_x^{*\sim})$.

Thus, if $A\in\mathcal{A}_x$, the datum of a correlation amounts to the
datum of an element of $H(A)$, i.e. of an isomorphism
\deleted{\unsure{of}} between the $T_x$-torsors
$\unsure{\mathbb{S}}_x^{\sim}(A)$ and
$\unsure{\mathbb{D}}_x^{*\sim}(\unsure{\mathrm{Ep}|E_p}(A))$.
This identification is compatible with the inclusion relations
between arcs $A\subset A'$.
\deleted{\unsure{D}} The correlations between
$\unsure{\mathbb{S}}_x^{\unsure{*|\sim}}$ and
$\unsure{\mathbb{D}}_x^{*\sim}$ form a torsor \unsure{under|over} $T_x$.
\deleted{For $\unsure{\mathbb{S}}_x^{\sim}$ fixed}

Let $\underline{\mathrm{Corr}}_x(\unsure{\mathfrak{X}|X})
=\underline{\mathrm{Corr}}_x(\deleted{\unsure{\mathbb{S}}}\unsure{X})$
be the category of \deleted{\unsure{pairs}}\unsure{triples}
$(\widetilde{\unsure{\mathbb{S}}}_x,
\widetilde{\unsure{\mathbb{D}}}_x^*,h)$ of correlated universal
coverings of $\unsure{\mathbb{S}}_x$ and of $\unsure{\mathbb{D}}_x^*$,
we have \deleted{\unsure{equi|equivalences}} functors
\[
\diagram{\begin{tikzcd}[ampersand replacement=\&, column sep=large,row sep=0.2ex]
\& \unsure{\mathbb{S}}_x^{\unsure{*\sim|\sim}} \\
(\widetilde{\unsure{\mathbb{S}}}_x,\unsure{\mathbb{D}}_x^{*\sim},h)
\arrow[ru] \arrow[rd] \& \\
\& \unsure{\mathbb{D}}_x^{*\sim}
\end{tikzcd}}
\]
\[
(117)\quad
\diagram{\begin{tikzcd}[ampersand replacement=\&, column sep=large,row sep=large]
\& \underline{\mathrm{Corr}}_x(X) \arrow[ld,"\approx"'] \arrow[rd,"\approx"] \& \\
\underline{\unsure{\mathrm{Revuni}}}(\unsure{\mathbb{S}}_x)
\& \& \underline{\unsure{\mathrm{Revuni}}}
(\deleted{\unsure{\mathbb{S}}}\unsure{\mathbb{D}}_x^*)
\end{tikzcd}}
\;,
\]
\unsure{which} are equivalences of categories.

\margin{268}

\margin{\begin{tabular}{c}
268\\
\texttt{https://grothendieck.umontpellier.fr}\\
Université de Montpellier, UMR CNRS 5149, Case courrier 051, Place Eugène Bataillon\\
34095 Montpellier cedex 5 - France
\end{tabular}}

\scanpage{270}

All these constructions have been carried out intrinsically carefully
enough to \unsure{pass} to isomorphisms of germs of differentiable
surfaces
$\diagram{\begin{tikzcd}[ampersand replacement=\&, column sep=large]
(X,x) \arrow[r,"\substack{u\\\sim}"] \& (X',x')
\end{tikzcd}}$.
Suppose chosen, in addition to
$(\unsure{\mathbb{S}}_x^{\sim},\unsure{\mathbb{D}}_x^{*\sim},h_x)$
\unsure{relative|relative} to $(X,x)$, also
$(\unsure{\mathbb{S}}_{x'}^{\sim},\unsure{a}\,
\unsure{\mathbb{D}}_{x'}^{*\sim},h'_{x'})$ \unsure{relative|relative} to $(X',x')$.
\unsure{Now} the isomorphism of the germs being given, it therefore induces
\[
(118)\quad\left\{
\diagram{\begin{array}{lll}
u_{\unsure{\mathbb{S}}}:&
\unsure{\mathbb{S}}_x\simeq\unsure{\mathbb{S}}_{x'}
&\text{isomorphism of topological spaces}\\[0.7em]
u_{\unsure{\mathbb{D}}^*}:&
\unsure{\mathbb{D}}_x^*\xrightarrow{\sim}\unsure{\mathbb{D}}_{x'}^*
&\text{isomorphism of germs of spaces}
\end{array}}
\right.
\]
\[
\diagram{\begin{array}{lll}
u_T:&T_x\simeq T_{x'}
&\begin{array}{l}
\text{isomorphism of the }\mathbb{Z}\text{-modules }\inserted{\text{topological}}\\
\unsure{\text{free}}\text{ of rank }1
\end{array}
\end{array}}
\]
Let us then consider the two sets
\[
(119)\quad\left\{
\begin{array}{l}
F(u)=\text{set of }\unsure{\text{isomorphisms}}\quad
\unsure{\mathbb{S}}_x^{\sim}\simeq\unsure{\mathbb{S}}_{x'}^{\sim}
\text{ over }\inserted{u_{\unsure{\mathbb{S}}}}
\\[0.6em]
\hfill\deleted{\unsure{u_{\mathbb{S}}}}
\\[0.6em]
G(u)=\text{set of isomorphisms}\quad
\unsure{\mathbb{D}}_x^{*\sim}\simeq\unsure{\mathbb{D}}_{x'}^{*\sim}
\text{ over }\inserted{u_{\unsure{\mathbb{D}}^*}},
\end{array}\right.
\]
they are both $T_x$-torsors.
This being so, I say that we have a canonical isomorphism of $T_x$-torsors
\[
(120)\quad
\diagram{\begin{tikzcd}[ampersand replacement=\&, column sep=large]
F(u) \arrow[r,"\sim"] \& G(u)
\end{tikzcd}}
\]
--- i.e. ``it comes to the $=$ to lift $u_{\unsure{\mathbb{S}}}$ to
$u_{\widetilde{\unsure{\mathbb{S}}}}:
\unsure{\mathbb{S}}_x^{\sim}\simeq S_{x'}^{\sim}$,
or to lift $u_{\unsure{\mathbb{D}}^*}$ to
\deleted{$\unsure{u_{\widetilde{\mathbb{S}}}}$}
$u_{\unsure{\mathbb{D}}^{*\sim}}:
\unsure{\mathbb{D}}_x^{*\sim}\unsure{\longrightarrow|-}
\unsure{\mathbb{D}}_{x'}^{*\sim}$''.
To see this, we note that by
\deleted{functoriality of the \unsure{constructions}}
\inserted{transport of structure} the

\margin{269}

\margin{\begin{tabular}{c}
269\\
\texttt{https://grothendieck.umontpellier.fr}\\
Université de Montpellier, UMR CNRS 5149, Case courrier 051, Place Eugène Bataillon\\
34095 Montpellier cedex 5 - France
\end{tabular}}

\scanpage{271}
\mspage{783}

inverse images under $u_{\unsure{\mathbb{S}|\mathcal{S}}}$ and under $u_{\unsure{\mathbb{D}|D}^{*}}$ of
$\unsure{\mathbb{S}|\mathcal{S}}_{x'}^{\sim}$ and of
$\unsure{\mathbb{D}|D}_{x'}^{*\sim}$,
\unsure{that is|that is} $\unsure{\mathbb{S}|\mathcal{S}}_x^{\prime\sim}$, $\unsure{\mathbb{D}|D}_x^{*\prime\sim}$,
are equipped with a correlation, which I denote by $h'_x$.
This posited, \deleted{\unsure{the existence} of an isomorphism
$F(u)\simeq G(u)$} canonical isomorphisms of $T_x$-torsors
\[
\begin{array}{rl}
(121)&F(u)\simeq\operatorname{Isom}_{\unsure{\mathbb{S}|\mathcal{S}}_x}
  (\unsure{\mathbb{S}|\mathcal{S}}_x^{\sim},\unsure{\mathbb{S}|\mathcal{S}}_x^{\prime\sim})\\[3pt]
&G(u)\simeq\operatorname{Isom}_{\unsure{\mathbb{D}|D}_x^{*}}
  (\unsure{\mathbb{D}|D}_x^{*\sim},\unsure{\mathbb{D}|D}_x^{*\prime\sim})
\end{array}
\]
Now by the diagram of equivalences (117) we do have isomorphisms of
$T_x$-torsors
\[
(122)\quad
\diagram{\begin{tikzcd}[ampersand replacement=\&, column sep=small,row sep=large]
\& {\operatorname{Isom}((\unsure{\mathbb{S}|\mathcal{S}}_x^{\sim},\unsure{\mathbb{D}|D}_x^{*\sim},h_x),
(\unsure{\mathbb{S}|\mathcal{S}}_x^{\prime\sim},\unsure{\mathbb{D}|D}_x^{*\prime\sim},h'_x))}
\arrow[dl,"\sim"']\arrow[dr,"\sim"] \& \\
{F(u)=\operatorname{Isom}(\unsure{\mathbb{S}|\mathcal{S}}_x^{\sim},\unsure{\mathbb{S}|\mathcal{S}}_x^{\prime\sim}\unsure{]|)}}
\&\&{G(u)=\unsure{\operatorname{Isom}}(\unsure{\mathbb{D}|D}_x^{*\sim},\unsure{\mathbb{D}|D}_x^{*\prime\sim})}
\end{tikzcd}}
\quad ,
\]
whence the canonical bijection \unsure{sought} (120) ---
\deleted{\unsure{of the}} which gives in particular (96).

\deleted{\unsure{The}} The canonical bijection (120) has the following
$\pm$ obvious property, which characterizes it:
\[
(123)\quad\left\{\vcenter{\hbox{\begin{minipage}{0.82\linewidth}
\deleted{\unsure{We}} Let $u_{\unsure{\mathbb{S}|\mathcal{S}}^{\sim}}$ and
$u_{\unsure{\mathbb{D}|D}^{*\sim}}$ be
\deleted{\unsure{extensions}}\inserted{liftings} of $u$
\deleted{\unsure{$u_{\mathbb{S}}$}}\deleted{\unsure{associated}}
$\unsure{\mathbb{S}|\mathcal{S}}_x^{\sim}\simeq\unsure{\mathbb{S}|\mathcal{S}}_{x'}^{\sim}$,
$\unsure{\mathbb{D}|D}_x^{*\sim}\longrightarrow\unsure{\mathbb{D}|D}_{x'}^{*\sim}$,
associated by (120), let $A\in\mathcal{A}_x$,
$A'\in\mathcal{A}'_x$ its image under
$u_{\unsure{\mathbb{S}|\mathcal{S}}}$, let $\alpha$ be a section of $\unsure{\mathbb{S}|\mathcal{S}}_x^{\sim}$ over $A$,
$\beta$ the section of $\unsure{\mathbb{D}|D}_x^{*\sim}$ over
$\unsure{\mathrm{Ep}|E_p}(A)$ \unsure{associated|associated} by the correlation
$h_x$, let $\alpha'$ be the section of $\unsure{\mathbb{S}|\mathcal{S}}_{x'}^{\sim}$ over $A'$
deduced from $\alpha$ by applying $u_{\unsure{\mathbb{S}|\mathcal{S}}^{\sim}}$,
$\beta'$ the section of $\unsure{\mathbb{D}|D}_{x'}^{*\sim}$ over
$\mathrm{Ep}(A')$ deduced from $\beta$ by applying
$u_{\unsure{\mathbb{D}|D}^{*\sim}}$, then $\alpha'$ and $\beta'$ are
\unsure{associated} by the correlation $h'$.
\end{minipage}}}\right.
\]
\margin{270}

\margin{\begin{tabular}{c}
270\\
\texttt{https://grothendieck.umontpellier.fr}\\
Université de Montpellier, UMR CNRS 5149, Case courrier 051, Place Eugène Bataillon\\
34095 Montpellier cedex 5 - France
\end{tabular}}

\scanpage{272}

This \deleted{\unsure{characterizes}} \unsure{(}for $A,\alpha$ fixed)
characterizes the map \deleted{bijection} (120), because for
$u_{\unsure{\mathbb{S}|\mathcal{S}}^{\sim}}$ fixed, it determines $u_{\unsure{\mathbb{D}|D}^{*\sim}}$ as follows:
we have $\beta$ by $h_x$, $\alpha'$ \unsure{deduced|deduced} from $\alpha$ by
$u_{\unsure{\mathbb{S}|\mathcal{S}}}$, $\beta'$ \unsure{deduced|deduced} from $\alpha'$ by $h_{x'}$,
now $\beta'$ must be equal to the section deduced from $\beta$ by applying
$u_{\unsure{\mathbb{D}|D}^{*\sim}}:\unsure{\mathbb{D}|D}_x^{*\sim}
\unsure{\longrightarrow|-}\unsure{\mathbb{D}|D}_{x'}^{*\sim}$,
and this indeed determines $u_{\unsure{\mathbb{D}|D}^{*\sim}}$ uniquely ...

\uline{N.B.} We can use this construction, for $A,\alpha$ fixed,
to define the map (120), we find that it is a
$T_x$-morphism (hence an isomorphism) and,
\deleted{\unsure{that it|since it}}\inserted{at the \unsure{same|=}}
time, that it does not depend on the choice of $\alpha$.
But a priori it depends \unsure{perhaps} on the choice of $A$, and the point which
requires verification is that this bijection in fact does not depend on the choice
of $A$. We know indeed that for $A\subset A'$, it is the
\unsure{same|=} for $A$ and $A'\unsure{;|,}$
whence at once the result, for the equivalence relation generated in
$\mathcal{A}_x$ by the inclusion relation is the indiscrete relation .... .

\margin{271}

\margin{\begin{tabular}{c}
271\\
\texttt{https://grothendieck.umontpellier.fr}\\
Université de Montpellier, UMR CNRS 5149, Case courrier 051, Place Eugène Bataillon\\
34095 Montpellier cedex 5 - France
\end{tabular}}